\section{Project Overview}
\label{section:intro}
\subsection{Profile}
\paragraph{Track: 3. Application.}
We study evidence selection for sentence-level hallucination classification using a fixed small language model.

\paragraph{Abstract and project goal.}
LLM-generated answers can contain statements that contradict or lack support in reference documents.
This project investigates whether selecting targeted evidence improves sentence-level hallucination classification on the public ANAH dataset.
We compare question-and-sentence-only input, leading-document truncation, BM25 retrieval, and BM25 followed by neural reranking.
All conditions share the same classifier, label definitions, prompt, and decoding settings, while the evidence-bearing conditions share a maximum evidence budget.
We measure four-class macro-F1, class-specific performance, processing cost, and evidence-selection failures.
The intended outcome is a reproducible analysis of when retrieval and reranking help classification and whether their benefits justify their additional cost.

\paragraph{TL;DR.}
With a fixed classifier and evidence budget, compare document truncation, BM25 retrieval, and neural reranking for sentence-level hallucination classification.

\subsection{Motivation}
A fluent answer is not necessarily faithful to its sources. A classifier reading only an answer sentence may rely on linguistic cues or parametric memory, whereas a classifier reading only the beginning of a document may miss decisive evidence elsewhere. Selecting useful evidence within a limited context is therefore a practical problem for reliable QA.
\begin{itemize}
\item \textbf{Challenge 1: Finding decisive evidence.} Lexical or semantic relevance alone does not establish whether a passage supports or contradicts a claim.
\item \textbf{Challenge 2: Retrieval failure versus unverifiability.} Evidence may exist in the full source but be absent from the selected context. The original dataset label remains the evaluation target; retrieval failure must not redefine it.
\item \textbf{Challenge 3: Fair comparisons and class imbalance.} Evidence selection changes the classifier's input distribution, and majority-class accuracy can hide minority-class failures. We control inference settings and report class-specific results.
\end{itemize}

\paragraph{Prior work and open question.}
SciFact established evidence-based scientific claim verification \citep{wadden2020}; BM25 is a foundational lexical retrieval method \citep{robertson2009}; and RAG demonstrates the value of combining retrieved documents with neural generation \citep{lewis2020}.
ANAH provides sentence-level factuality labels, reference fragments, and corrections for QA responses \citep{ji2024anah}. Its annotation pipeline already uses retrieval, so we do not claim to introduce retrieval-based hallucination checking.
Recent work includes iterative retrieval and verification in FIRE \citep{xie2025fire} and listwise reranking in jina-reranker-v3 \citep{wang2025jina}.
Our open question is whether these evidence-selection choices improve a fixed small classifier under the same evidence budget, which classes benefit, and what costs are incurred.

\paragraph{State-of-the-art reference and scope.}
We select jina-reranker-v3 as a modern, strong retrieval comparison based on its reported retrieval results, not as an established best classifier on ANAH.
FIRE is related verification work rather than a required implementation. Our contribution is a controlled application study; results under different models, taxonomies, or dataset splits will not be treated as directly comparable leaderboard scores.
